% =========================================================
% theorem2_proof_skeleton.tex
% Theorem 2 (Impossibility) — Proof Skeleton
% Project: PrecedentGuard / AAAI 2027 submission
% Created: 2026-06-30
% Status: SKELETON — fill in proof bodies marked [FILL]
% Decision Gates: Day 7 (Gate 1), Day 14 (Gate 2)
% =========================================================

\documentclass[11pt]{article}
\usepackage[margin=1in]{geometry}
\usepackage{amsmath, amssymb, amsthm, mathtools}
\usepackage{algorithm}
\usepackage[noend]{algpseudocode}
\usepackage{hyperref}
\usepackage{tikz}
\usepackage{xcolor}
\usepackage{enumitem}

\newcommand{\TODO}[1]{{\color{red}\textbf{[FILL: #1]}}}
\newcommand{\NOTE}[1]{{\color{blue}\textit{[Note: #1]}}}
\newcommand{\REF}[1]{{\color{teal}\textsf{[Ref: #1]}}}

\theoremstyle{plain}
\newtheorem{theorem}{Theorem}
\newtheorem{lemma}{Lemma}
\newtheorem{proposition}{Proposition}
\newtheorem{corollary}{Corollary}
\theoremstyle{definition}
\newtheorem{definition}{Definition}
\newtheorem{assumption}{Assumption}
\newtheorem{remark}{Remark}

\DeclareMathOperator{\TV}{TV}
\DeclareMathOperator{\KL}{KL}
\DeclareMathOperator{\negl}{negl}
\DeclareMathOperator{\Adv}{Adv}
\DeclareMathOperator{\Pa}{Pa}
\DeclareMathOperator{\Ch}{Ch}
\DeclareMathOperator{\do}{do}

\title{Theorem 2 (Impossibility) Proof Skeleton:\\
       \large No Double-Sided Bounded Do-Intervention without a Cryptographically-Isolated Trust Anchor}
\author{PrecedentGuard --- AAAI 2027 Submission Draft}
\date{2026-06-30}

\begin{document}
\maketitle

\tableofcontents
\newpage

% =========================================================
\section{Preliminaries and Notation}
% =========================================================

\subsection{Causal Influence Diagram for an Agent Trajectory}

\begin{definition}[Causal Influence Diagram, CID]\label{def:cid}
A \emph{Causal Influence Diagram} for an LLM agent trajectory is a tuple
$\mathcal{G} = (\mathcal{V}, \mathcal{E}, \mathcal{T})$ where
$\mathcal{V}$ is a finite set of nodes,
$\mathcal{E} \subseteq \mathcal{V}\times\mathcal{V}$ is a set of directed edges,
and $\mathcal{T}: \mathcal{V} \to \{\textsc{Intent}, \textsc{Memory}, \textsc{Tool}, \textsc{Obs}, \textsc{Decision}, \textsc{Utility}\}$
is a node-type assignment. We write $\Pa(v)$ for the parent set of node $v\in\mathcal{V}$
and $\Ch(v)$ for its child set.
\end{definition}

\REF{Everitt et al.\ 2021, Synthese; Pearl 2009 \S 1.2.}

\begin{remark}\label{rem:dnodes}
Within a single trajectory, the \emph{decision node} $D\in\mathcal{V}$ is the
guard's binary output $D \in \{\textsc{allow}, \textsc{block}\}$. The \emph{utility node}
$U$ encodes the safety-correctness label (1 if the underlying trajectory is
actually unsafe, 0 if safe).
\end{remark}

\subsection{Memory Augmentation and the Precedent Tuple}

\begin{definition}[Precedent Tuple]\label{def:precedent}
Each retrievable precedent is a five-tuple
$E_i = (M_i, G_i, T_i, L_i, C_i)$, where
$M_i \in \mathbb{R}^{d_M}$ is a dense semantic vector,
$G_i$ is a typed causal sub-graph,
$T_i$ is a timestamp,
$L_i \in \{0,1\}^*$ is a (possibly empty) cryptographically signed audit rule,
and $C_i \in \{0,1\}$ is a binary safety label.
We partition $E_i$ into two tiers:
the \emph{poisonable tier} $\Pi_{\mathrm{P}}(E_i) := \{M_i, G_i, T_i, C_i\}$
and the \emph{signed tier} $\Pi_{\mathrm{S}}(E_i) := \{L_i\}$.
\end{definition}

\subsection{Bounded Do-Intervention Scheme}

\begin{definition}[Bounded Do-Intervention Scheme]\label{def:bdis}
A \emph{bounded do-intervention scheme} is a tuple
$\Pi = (R, \Phi, \varepsilon)$, where
\begin{itemize}[noitemsep]
  \item $R: \mathcal{X} \times \mathcal{B} \to \mathcal{B}^K$ is a retrieval map
        from query trajectory and memory bank to a top-$K$ retrieved set;
  \item $\Phi: \mathcal{B}^K \times \mathbb{R}^{|D|} \to \mathbb{R}^{|D|}$
        is a logit perturbation map applied to a frozen guard
        $f_\theta: \mathcal{X} \to \mathbb{R}^{|D|}$;
  \item $\varepsilon > 0$ is a perturbation budget with the constraint
        $\|\Phi(\cdot,\cdot)\|_\infty \le \varepsilon$.
\end{itemize}
The guard's decision under $\Pi$ is
$D_\Pi(x) = \arg\max_k \bigl[f_\theta(x) + \Phi(R(x), f_\theta(x))\bigr]_k$.
\end{definition}

% =========================================================
\section{Threat Model}\label{sec:threat}
% =========================================================

\begin{definition}[Adversary $\mathcal{A}_{p, K}$]\label{def:adversary}
An \emph{intervention adversary} $\mathcal{A}_{p, K}$ is permitted to apply
\emph{do-operations} on at most $K$ causal-parent nodes of $D$
within a poisoning budget $p \in [0,1]$, where $p$ is the maximum fraction
of poisonable-tier content the adversary may corrupt.
\end{definition}

\begin{assumption}[No-trust-anchor setting]\label{ass:notrust}
The adversary has access only to the poisonable tier $\Pi_{\mathrm{P}}(E_i)$
and has no signing key. The defender, in the \emph{no-trust-anchor setting},
does not consult $\Pi_{\mathrm{S}}(E_i)$ when forming $\Phi$
(equivalently, $L_i$ does not influence the perturbation map).
\end{assumption}

\NOTE{This is the load-bearing condition for Theorem~2. It captures
the regime where the defender lacks a verifiable, unforgeable signal.}

\begin{definition}[FNR/FPR of $\Pi$]\label{def:fnrfpr}
Let $\mathcal{D}$ be the trajectory distribution. The false-negative and
false-positive rates of $\Pi$ are
\[
  \mathrm{FNR}(\Pi) = \Pr_{x\sim\mathcal{D}}\bigl[D_\Pi(x) = \textsc{allow} \mid U(x) = 1\bigr],
\quad
  \mathrm{FPR}(\Pi) = \Pr_{x\sim\mathcal{D}}\bigl[D_\Pi(x) = \textsc{block} \mid U(x) = 0\bigr].
\]
\end{definition}

% =========================================================
\section{Theorem 2 (Impossibility) --- Statement}
% =========================================================

\begin{theorem}[Impossibility of Double-Sided Bounded Do-Intervention]\label{thm:impossibility}
Let $\mathcal{G}$ be a Causal Influence Diagram satisfying Definition~\ref{def:cid},
and let $\Pi = (R, \Phi, \varepsilon)$ be any bounded do-intervention scheme
under Assumption~\ref{ass:notrust} (no-trust-anchor setting).
For any $\delta < 1/2$, no $\Pi$ can simultaneously satisfy
\[
  \mathrm{FNR}(\Pi) \le \delta \quad \text{and} \quad \mathrm{FPR}(\Pi) \le \delta
\]
against an adversary $\mathcal{A}_{p, K}$ with $p \cdot K \ge p_*$,
where $p_*$ is a constructive threshold defined in the proof.
\end{theorem}

\begin{remark}
The constant $p_*$ depends on $K$, the embedding dimension $d_M$,
and the structure of the retrieval map $R$. \TODO{Provide explicit
formula for $p_*$ once the construction in Lemma~\ref{lem:construct} is filled in.}
\end{remark}

% =========================================================
\section{Supporting Lemmas}
% =========================================================

\subsection{Lemma 1 --- Existence of Indistinguishable Intervention Pair}

\begin{lemma}[Adversarial Indistinguishable Construction]\label{lem:construct}
Under Assumption~\ref{ass:notrust}, there exist two trajectories $x_1, x_2$
and a do-intervention $\nu \in \mathcal{A}_{p, K}$ such that
\begin{enumerate}[label=(\roman*), noitemsep]
  \item $U(x_1) = 0$ (safe ground truth) and $U(x_2) = 1$ (unsafe ground truth);
  \item the retrieved top-$K$ sets satisfy
        $R(x_1) \overset{d}{=} R(\nu(x_2))$ as distributions
        on the retrieval space;
  \item the perturbation map $\Phi$ takes input distributions only,
        so $\Phi(R(x_1), \cdot) \overset{d}{=} \Phi(R(\nu(x_2)), \cdot)$.
\end{enumerate}
\end{lemma}

\begin{proof}[Proof sketch]
\TODO{Construct $x_2$ as an unsafe trajectory whose natural memory pool
$\mathcal{B}_2$ contains $K$ items semantically distant from $x_1$'s
natural retrieved set. Then let the adversary inject $K$ crafted items
into $\mathcal{B}_2$ such that retrieval similarity to $x_1$'s top-$K$
is maximized.}

\TODO{Argue that under the retrieval map's continuity (Lipschitz or
nearest-neighbor), the resulting retrieved set is identically distributed
to $R(x_1)$. Use the embedding-space Voronoi argument from \REF{Bhagoji et al.\ 2018, NeurIPS} (\S 3.1)
adapted to typed-set similarity.}

\TODO{Bound the required injection size: $K_{\mathrm{inject}} \le K$, which
holds whenever $p \cdot |\mathcal{B}_2| \ge K$, giving the threshold
$p_* = K / |\mathcal{B}_2|$.}
\end{proof}

\subsection{Lemma 2 --- Decision-Coincidence under Indistinguishability}

\begin{lemma}[Decision Coincidence]\label{lem:coincide}
Let $\Pi = (R, \Phi, \varepsilon)$ be any bounded do-intervention scheme
under Assumption~\ref{ass:notrust}. For the $x_1, x_2, \nu$ given by
Lemma~\ref{lem:construct},
\[
  D_\Pi(x_1) = D_\Pi(\nu(x_2))
  \quad\text{with probability 1.}
\]
\end{lemma}

\begin{proof}[Proof sketch]
\TODO{Since $R(x_1) \overset{d}{=} R(\nu(x_2))$ and $\Phi$ depends on $R$ only
(no signed $L_i$ access by Assumption~\ref{ass:notrust}),
$\Phi(R(x_1), \cdot) \overset{d}{=} \Phi(R(\nu(x_2)), \cdot)$.}

\TODO{The frozen guard $f_\theta$ is a deterministic function of the
trajectory's text representation. Argue that on $x_1$ and $\nu(x_2)$
the text representation of the retrieved memory chunk fed back into
the guard input is identical (or has identical distribution).}

\TODO{Therefore $\arg\max_k [f_\theta + \Phi]_k$ yields the same decision.}
\end{proof}

\begin{remark}
Lemma~\ref{lem:coincide} is the load-bearing step. It transforms a
\emph{statistical} indistinguishability (Lemma~\ref{lem:construct})
into a \emph{decision-level} coincidence.
\end{remark}

\subsection{Lemma 3 --- Lower Bound on the Maximum of FNR and FPR}

\begin{lemma}[Maximum-Error Lower Bound]\label{lem:lowerbound}
Under the conditions of Lemmas~\ref{lem:construct}--\ref{lem:coincide},
\[
  \max\bigl(\mathrm{FNR}(\Pi), \mathrm{FPR}(\Pi)\bigr) \ge \frac{1}{2}.
\]
\end{lemma}

\begin{proof}[Proof sketch]
\TODO{By Lemma~\ref{lem:coincide}, $\Pi$ outputs the same decision on
$x_1$ (safe) and $\nu(x_2)$ (unsafe). Without loss of generality,
denote this common decision $d^\star$.}

\TODO{If $d^\star = \textsc{allow}$, then on $\nu(x_2)$ we have
$D_\Pi = \textsc{allow}$ while $U = 1$ --- a false negative event.}

\TODO{If $d^\star = \textsc{block}$, then on $x_1$ we have
$D_\Pi = \textsc{block}$ while $U = 0$ --- a false positive event.}

\TODO{Mix over the joint distribution of $(x_1, x_2)$ in the
adversary's strategy: at least one of FNR or FPR exceeds $1/2$
because the constructed pair contributes equally to both rates
under a balanced reduction.}

\REF{This argument follows the Bayes-error / minimax structure of
\REF{Tsipras et al.\ 2019, ICLR} \S 3 and the indistinguishability
arguments of \REF{Cohen, Rosenfeld \& Kolter 2019, ICML} \S 3.}
\end{proof}

\subsection{Lemma 4 --- Trust Anchor Restores Distinguishability}

\begin{lemma}[Trust Anchor Distinguishability]\label{lem:anchor}
Assume the signature scheme used to produce $L_i$ is
existentially unforgeable under chosen-message attack (EUF-CMA),
with security parameter $\lambda$.
Then, with access to $\Pi_{\mathrm{S}}(E_i) = \{L_i\}$, the indistinguishability
of Lemma~\ref{lem:construct} no longer holds: there is a verifier
$V$ such that
\[
  \Pr\bigl[V(R(x_1)) = 1 \wedge V(R(\nu(x_2))) = 0\bigr] \ge 1 - \negl(\lambda),
\]
where $\negl(\lambda)$ is a function vanishing faster than any inverse polynomial.
\end{lemma}

\begin{proof}[Proof sketch]
\TODO{The genuine precedent $E_1$ retrieved for $x_1$ carries a valid
$L_1$ signed by the trusted authority. The adversary, lacking the
signing key, cannot produce a $L_2'$ such that $\mathrm{Verify}(L_2', m) = 1$
for the message $m$ corresponding to the injected memory --- except
with probability $\negl(\lambda)$ by EUF-CMA.}

\TODO{Therefore the verifier $V$ accepts $R(x_1)$ but rejects
$R(\nu(x_2))$ with overwhelming probability, breaking the
decision-coincidence chain of Lemma~\ref{lem:coincide}.}

\REF{Standard EUF-CMA argument; see Katz \& Lindell, \textit{Introduction to
Modern Cryptography} (3rd ed.), \S 12.4.}
\end{proof}

% =========================================================
\section{Proof of Theorem 2}
% =========================================================

\begin{proof}[Proof of Theorem~\ref{thm:impossibility}]
The proof proceeds in four steps, instantiating the lemmas above.

\paragraph{Step 1 --- Setup.}
Fix the no-trust-anchor setting of Assumption~\ref{ass:notrust}.
Let $\Pi = (R, \Phi, \varepsilon)$ be an arbitrary bounded do-intervention scheme.

\paragraph{Step 2 --- Construction.}
By Lemma~\ref{lem:construct}, there exist $x_1$ (safe), $x_2$ (unsafe),
and an adversarial intervention $\nu \in \mathcal{A}_{p, K}$ with $p \ge p_*$
producing the indistinguishable retrieval distributions.

\paragraph{Step 3 --- Decision Coincidence.}
By Lemma~\ref{lem:coincide}, $D_\Pi(x_1) = D_\Pi(\nu(x_2))$ with probability 1.

\paragraph{Step 4 --- Lower Bound.}
By Lemma~\ref{lem:lowerbound}, $\max(\mathrm{FNR}, \mathrm{FPR}) \ge 1/2$.
Since this contradicts any joint bound of the form
$\mathrm{FNR} \le \delta \wedge \mathrm{FPR} \le \delta$ for $\delta < 1/2$,
the desired impossibility follows.

\paragraph{Contrast with Trust-Anchor Setting.}
Under access to $\Pi_{\mathrm{S}}(E_i)$, Lemma~\ref{lem:anchor} restores
distinguishability with overwhelming probability, and the double-sided
bound of Theorem~1 becomes achievable. This establishes the trust anchor
as \emph{necessary}, not merely sufficient.
\qedhere
\end{proof}

% =========================================================
\section{Corollaries (Optional, for Section 4 of Paper)}
% =========================================================

\begin{corollary}[Single-Sided Possibility under Stronger Assumption]\label{cor:single}
\TODO{State that under bounded statistical distance (e.g.\ Lipschitz retrieval),
a \emph{single-sided} bound (FNR only, as in SMSR) remains achievable without
a trust anchor.}
\end{corollary}

\begin{corollary}[Threshold $p_*$ Dependence]\label{cor:threshold}
\TODO{Express $p_*$ as a function of $K$, $|\mathcal{B}|$, and embedding-space
geometry. Useful for empirical $p^*$ matching in \S 4.4.}
\end{corollary}

\begin{corollary}[Equivalence with the Static-Guard Limit]\label{cor:static}
\TODO{Argue that for $K = 0$ (no memory retrieval, i.e.\ a static guard),
the impossibility reduces to the standard Bayes-error lower bound.}
\end{corollary}

% =========================================================
\section{Day-by-Day Task Checklist (14 Days)}
% =========================================================

\paragraph{Week 1 --- Construction and Indistinguishability}

\begin{enumerate}[label=\textbf{Day \arabic*.}, leftmargin=*]
  \item \textbf{Notation lockdown.} Finalize Definitions 1--4 with PrecedentGuard v2
        notation. Verify no clash with existing paper.
  \item \textbf{Threat model formalization.} Write Assumption~\ref{ass:notrust}
        and Definition~\ref{def:adversary} in full. Identify all variables.
  \item \textbf{Lemma 1 construction.} Construct $x_1, x_2, \nu$ on a toy CID
        (3 memory nodes, 2 tools). Verify retrieval similarity by hand.
  \item \textbf{Lemma 1 generalization.} Lift the toy construction to the
        general CID. Establish $p_* = K / |\mathcal{B}_2|$.
  \item \textbf{Lemma 1 LaTeX.} Write full proof sketch with embedding-space
        Voronoi argument. \REF{Bhagoji 2018 \S 3.1.}
  \item \textbf{Lemma 2 setup.} Argue $\Phi$'s dependence is only on $R(\cdot)$.
        Show this implies $\Phi$-output coincidence.
  \item \textbf{Gate 1 review.} Self-audit Lemmas 1--2.
        \textbf{Decision:} are constructions sound? If not, pivot to Plan B.
        \\\textbf{Stop-loss:} if construction fails after Day 7, drop Theorem~2 and finalize Theorem~1 + corollaries only.
\end{enumerate}

\paragraph{Week 2 --- Lower Bound, Cryptographic Step, and Verification}

\begin{enumerate}[label=\textbf{Day \arabic*.}, leftmargin=*, start=8]
  \item \textbf{Lemma 3 derivation.} Write the FNR/FPR lower-bound argument.
        Reference \REF{Tsipras 2019 \S 3} and \REF{Cohen-Rosenfeld-Kolter 2019 \S 3}.
  \item \textbf{Lemma 4 cryptography.} Invoke EUF-CMA for Ed25519
        (PrecedentGuard v2's signature scheme). Cite Katz \& Lindell \S 12.4.
  \item \textbf{Theorem 2 main proof.} Assemble Lemmas 1--4 into the four-step
        proof. Verify each link in the chain.
  \item \textbf{Corollary draft.} Write Corollaries 1--3.
  \item \textbf{Internal review.} Hand the draft to one math-strong colleague
        (or LLM self-review with adversarial prompt). Mark any unclear step.
  \item \textbf{Revision.} Fix all flagged steps. Tighten constants.
  \item \textbf{Gate 2 review.} Final audit.
        \textbf{Decision:} is the proof airtight?
        \\\textbf{Outcome A:} Theorem~2 holds $\Rightarrow$ Plan A
        (proceed to \S 3 Method with double-theorem narrative).
        \\\textbf{Outcome B:} Theorem~2 has fixable gap $\Rightarrow$
        extend 1 week or downgrade to a weaker version
        (e.g.\ ``one-sided impossibility'').
        \\\textbf{Outcome C:} Theorem~2 fundamentally fails $\Rightarrow$
        commit to Plan B (Theorem~1 + 3 corollaries only).
\end{enumerate}

% =========================================================
\section{Plan B Fallback (If Theorem 2 Fails)}
% =========================================================

If Theorem~2 cannot be established by Day 14, the paper's contribution
structure must be revised:

\begin{itemize}[noitemsep]
  \item \textbf{C1} stays --- CID formalization is independent of Theorem~2.
  \item \textbf{C2} weakens to \emph{Theorem~1 + 3 corollaries only}.
  \item \textbf{C3} stays --- BDIC algorithm and empirical validation are unaffected.
  \item \textbf{\S 2.2 SMSR critique} must change from
        ``trust isolation is \emph{necessary}'' to
        ``trust isolation enables a \emph{stronger} guarantee''.
  \item \textbf{\S 3 Method} must motivate the trust anchor as an
        engineering choice rather than a theoretical necessity.
  \item \textbf{Target venue} may shift from AAAI 2027 to NeurIPS 2027
        for an extra 3 months of polishing.
\end{itemize}

\textbf{Verdict path priority:} Plan A (Theorem 2 holds) $\succ$
Plan B-strong (Theorem 2-weak: only single-direction impossibility holds) $\succ$
Plan B (Theorem 2 fails entirely).

% =========================================================
\section{Inline References (Minimal Pre-Reading)}
% =========================================================

The proof draws on the following four results. \textbf{Read only the cited section,
not the full paper}, unless intuition fails.

\begin{enumerate}[noitemsep]
  \item Madry, A. et al.\ (2018). \emph{Towards Deep Learning Models Resistant to
        Adversarial Attacks}. ICLR. \textbf{\S 2 only} --- adversarial budget formalization.
  \item Cohen, J., Rosenfeld, E., \& Kolter, J.Z.\ (2019). \emph{Certified Adversarial
        Robustness via Randomized Smoothing}. ICML. \textbf{\S 3 only} ---
        the certification-vs-attack indistinguishability structure.
  \item Tsipras, D. et al.\ (2019). \emph{Robustness May Be at Odds with Accuracy}.
        ICLR. \textbf{\S 3 only} --- impossibility-style argument with Bayes-error bound.
  \item Bhagoji, A.N., Cullina, D., \& Mittal, P.\ (2018). \emph{Lower Bounds on Adversarial
        Robustness from Optimal Transport}. NeurIPS. \textbf{\S 3.1 only} ---
        embedding-space construction technique.
\end{enumerate}

\REF{Optional deeper background:} Everitt, T.\ et al.\ (2021), CID survey;
Pearl, J.\ (2009), \textit{Causality}, \S 3.4 do-operator semantics;
Katz, J.\ \& Lindell, Y.\ (2020), \textit{Introduction to Modern Cryptography}, \S 12 EUF-CMA.

% =========================================================
\section{Self-Review Checklist (Before Gate 2)}
% =========================================================

\begin{itemize}[label=$\Box$, noitemsep]
  \item Definitions 1--4 are notation-compatible with PrecedentGuard v2 \S 2--3.
  \item Assumption~\ref{ass:notrust} clearly excludes only the signed tier $L_i$.
  \item Lemma 1's construction is concrete enough to compile a toy example in Python.
  \item Lemma 2's argument explicitly addresses both deterministic and randomized
        guards.
  \item Lemma 3's $1/2$ bound is matched by a constructive adversary, not just stated.
  \item Lemma 4's EUF-CMA invocation is rigorous (security parameter and reduction
        sketch present).
  \item Theorem 2's contradiction is sharp: $\max(\mathrm{FNR}, \mathrm{FPR}) \ge 1/2$
        cannot be relaxed.
  \item Corollary 1 is consistent with SMSR's published single-sided bound
        (sanity check).
  \item All three inline reference \S sections have been re-read.
  \item One colleague has independently verified Lemmas 1--4.
\end{itemize}

\end{document}
